\chapter{Panoramica del sistema}
\label{ch:sistema}

\begin{semplice}
SYMPHONY V2 è una radio UHF per un CubeSat. Il microcontrollore (MCU) STM32L562 comanda la radio Si4464, che trasmette e riceve a \SI{435}{\mega\hertz}. In trasmissione il segnale viene amplificato da un amplificatore di potenza (PA). In ricezione passa da un amplificatore a basso rumore (LNA). Un commutatore RF (switch) collega l'antenna a uno dei due rami. Intorno ci sono una memoria (FRAM), un orologio (RTC), un watchdog, due sensori di temperatura e le interfacce verso il satellite (CAN, I2C) e verso il PC (USB).
\end{semplice}

\section{Schema a blocchi}
Lo schema è diviso in tre figure, per non affollare il disegno:
\begin{itemize}
\item figura~\ref{fig:blocchi}: l'MCU e tutto quello che le è collegato;
\item figura~\ref{fig:catena-rf}: il percorso del segnale radio;
\item figura~\ref{fig:alim-blocchi}: da dove arriva l'alimentazione.
\end{itemize}
Sulle frecce è scritta \textbf{la periferica dell'MCU} (per esempio FDCAN1), non il numero del pin. I pin di ogni periferica sono nella mappa dei pin (capitolo~\ref{ch:pin}) e nelle tabelle del capitolo~\ref{ch:periferiche}. Tutti i collegamenti sono presi dalla netlist dello schematico \fonte{SCH}.

\begin{figure}[htbp]
\centering
\begin{tikzpicture}[x=1mm,y=1mm,
   per/.style={blk,minimum width=34mm,minimum height=11mm},
   lab/.style={font=\sffamily\scriptsize,inner sep=1.5pt},
   bus/.style={{Stealth[length=2mm]}-{Stealth[length=2mm]},thick,sbblue2},
   uni/.style={-{Stealth[length=2mm]},thick,sbblue2},
   ctl/.style={-{Stealth[length=2mm]},thick,sbgreen}]
% MCU
\node[blk,minimum width=36mm,minimum height=104mm,font=\sffamily\small] (mcu) at (0,0)
  {\textbf{U301}\\\textbf{STM32L562}\\[3pt]microcontrollore\\[2pt]\footnotesize quarzo \SI{24}{\mega\hertz}};
% ---- sinistra
\node[per] (fram) at (-62,40) {U401 FM25V10\\memoria FRAM};
\node[per] (rtc)  at (-62,20) {U404 DS3231M\\orologio (RTC)};
\node[per] (wdt)  at (-62,0)  {U403 STWD100\\watchdog esterno};
\node[per] (adc)  at (-62,-20){IC1, IC5 TLV9051\\2 sensori PT1000};
\node[per] (swd)  at (-62,-40){J1 TC2030\\debugger};
\draw[bus] (fram.east) -- node[lab,above]{SPI2} (fram.east-|mcu.west);
\draw[bus] ($(rtc.east)+(0,2)$) -- node[lab,above]{I2C2} ($(rtc.east-|mcu.west)+(0,2)$);
\draw[uni] ($(rtc.east)+(0,-2)$) -- node[lab,below]{GPIO: avviso} ($(rtc.east-|mcu.west)+(0,-2)$);
\draw[uni] ($(wdt.east-|mcu.west)+(0,2)$) -- node[lab,above]{GPIO: impulso} ($(wdt.east)+(0,2)$);
\draw[ctl] ($(wdt.east)+(0,-2)$) -- node[lab,below,text=sbgreen]{reset (NRST)} ($(wdt.east-|mcu.west)+(0,-2)$);
\draw[uni] (adc.east) -- node[lab,above]{ADC1} (adc.east-|mcu.west);
\draw[bus] (swd.east) -- node[lab,above]{SWD} (swd.east-|mcu.west);
% ---- destra
\node[per] (si)  at (62,45)  {IC204 Si4464\\radio};
\node[per,blko] (sup) at (62,27)  {IC205 TPS3825\\supervisore radio};
\node[per,blko] (cmd) at (62,9)   {Q701 + U701\\switch, PA, LNA};
\node[per] (ch)  at (62,-9)  {U402 CH340E\\verso USB-C (J15)};
\node[per] (can) at (62,-27) {IC501 TCAN3414\\verso J202 e dock};
\node[per] (i2c) at (62,-45) {bus I2C esterno\\J203 e dock J708};
\draw[bus] ($(si.west-|mcu.east)+(0,2)$) -- node[lab,above]{SPI1} ($(si.west)+(0,2)$);
\draw[uni] ($(si.west)+(0,-2)$) -- node[lab,below]{EXTI (interrupt)} ($(si.west-|mcu.east)+(0,-2)$);
\draw[uni] (sup.west-|mcu.east) -- node[lab,above]{GPIO: reset radio} (sup.west);
\draw[uni] (cmd.west-|mcu.east) -- node[lab,above]{GPIO: TX/RX} (cmd.west);
\draw[bus] (ch.west-|mcu.east) -- node[lab,above]{USART1} (ch.west);
\draw[bus] (can.west-|mcu.east) -- node[lab,above]{FDCAN1} (can.west);
\draw[bus] (i2c.west-|mcu.east) -- node[lab,above]{I2C1} (i2c.west);
\draw[ctl] (sup.north) -- node[lab,right,text=sbgreen]{SDN} (si.south);
\end{tikzpicture}
\caption{L'MCU e le sue periferiche. Sulle frecce c'è la periferica dell'MCU usata. Blu: dati (la punta indica il verso, due punte = nei due versi). Verde: segnale che agisce da solo, senza passare dal firmware. Arancio: blocchi che proteggono la radio.}
\label{fig:blocchi}
\end{figure}

\begin{nota}[come leggere la figura~\ref{fig:blocchi}]
\begin{itemize}
\item \textbf{GPIO}: pin usato come semplice ingresso o uscita digitale, non come periferica. Dopo i due punti c'è a che cosa serve.
\item \textbf{GPIO: avviso}: il DS3231M abbassa questo pin quando la sua alimentazione cade (paragrafo~\ref{sec:rtc-rst}).
\item \textbf{GPIO: impulso}: il firmware deve cambiarne il livello almeno ogni secondo, altrimenti il watchdog resetta l'MCU tramite \textbf{NRST} (capitolo~\ref{ch:supervisione}).
\item \textbf{EXTI}: linea di interrupt esterno. Il Si4464 abbassa il pin quando ha un evento da segnalare (pacchetto inviato, pacchetto ricevuto\ldots).
\item \textbf{SDN}: il supervisore IC205 spegne la radio se il 3V3 scende troppo o se il firmware lo chiede con \enquote{GPIO: reset radio} (paragrafo~\ref{sec:sdn}).
\item \textbf{GPIO: TX/RX}: un solo pin decide se la scheda trasmette o riceve. Comanda insieme switch, PA e LNA (figura~\ref{fig:catena-rf} e paragrafo~\ref{sec:rf-arch}).
\end{itemize}
\end{nota}

\begin{figure}[htbp]
\centering
\begin{tikzpicture}[x=1mm,y=1mm,
   rb/.style={blk,minimum width=22mm,minimum height=11mm},
   rf/.style={-{Stealth[length=2mm]},very thick,sborange},
   lab/.style={font=\sffamily\scriptsize,inner sep=1.5pt},
   ctl/.style={-{Stealth[length=2mm]},thick,sbgreen},
   cl/.style={font=\sffamily\scriptsize,text=sbgreen,align=center}]
\node[blk,minimum width=20mm,minimum height=34mm] (si) at (0,-9) {IC204\\Si4464\\radio};
\node[blk,minimum width=20mm] (tcxo) at (0,22) {Y201 TCXO\\\SI{32}{\mega\hertz}};
\draw[-{Stealth[length=2mm]},thick,sbblue2] (tcxo) -- node[lab,right]{clock} (si);
% TX
\node[rb] (att) at (31,0)  {attenuatore\\\SI{10.8}{dB}};
\node[rb,blkr] (pa) at (58,0)  {IC2 QPA9510\\PA};
\node[rb] (lpf) at (85,0)  {FL201\\passa-basso};
% RX
\node[rb] (mat) at (31,-18) {adattamento\\RX};
\node[rb] (lna) at (58,-18) {Q201\\LNA};
\node[rb] (frx) at (85,-18) {filtro RX\\passa-basso};
\node[blko,minimum width=20mm,minimum height=34mm] (sw) at (114,-9) {IC202\\HMC349A\\switch};
\node[draw=sbgray,thick,circle,inner sep=1.5pt,font=\sffamily\scriptsize] (ant) at (139,-9) {MMCX};
\draw[rf] (si.east|-att) -- node[lab,above]{TX} (att);
\draw[rf] (att) -- (pa); \draw[rf] (pa) -- (lpf);
\draw[rf] (lpf) -- node[lab,above]{RF2} (lpf-|sw.west);
\draw[rf] (frx-|sw.west) -- node[lab,above]{RF1} (frx);
\draw[rf] (frx) -- (lna); \draw[rf] (lna) -- (mat);
\draw[rf] (mat) -- node[lab,above]{RX} (mat-|si.east);
\draw[very thick,sborange,{Stealth[length=2mm]}-{Stealth[length=2mm]}] (sw) -- node[lab,above]{RFC} (ant);
% comandi
\node[cl] (cpa) at (58,16) {comando TX/RX\\(accende il PA, tramite Q701)};
\node[cl] (csw) at (114,16) {comando TX/RX\\(sceglie il ramo)};
\node[cl] (cln) at (58,-34) {comando TX/RX\\(spegne l'LNA, tramite U701)};
\draw[ctl] (cpa) -- (pa); \draw[ctl] (csw) -- (sw);
\draw[ctl] (cln) -- (lna);
\end{tikzpicture}
\caption{Percorso del segnale radio. Sopra la trasmissione (TX), sotto la ricezione (RX). Lo switch collega l'antenna (connettore MMCX) a un ramo o all'altro: RF2 = trasmissione, RF1 = ricezione. In verde il comando TX/RX: è lo stesso pin dell'MCU per tutti e tre i blocchi.}
\label{fig:catena-rf}
\end{figure}

\begin{figure}[htbp]
\centering
\begin{tikzpicture}[x=1mm,y=1mm,
   pb/.style={blk,minimum width=30mm,minimum height=11mm},
   pwr/.style={-{Stealth[length=2mm]},very thick,sbred!80},
   lab/.style={font=\sffamily\scriptsize,inner sep=1.5pt},
   ctl/.style={-{Stealth[length=2mm]},thick,sbgreen}]
\node[pb] (dock) at (0,10)  {J708 dock\\(dal satellite)};
\node[pb] (usb)  at (0,-10) {J15 USB-C\\(dal PC)};
\node[pb] (mux)  at (46,0) {IC4 TPS2121\\scelta della sorgente};
\node[pb] (buck) at (94,0) {IC3 TPSM863257\\convertitore \SI{3.3}{\volt}};
\node[font=\sffamily\scriptsize,align=left,anchor=west] (load) at (120,0) {MCU, radio, PA,\\memorie, sensori};
\node[font=\sffamily\scriptsize,align=left,anchor=west] (swv) at (120,-16) {alimentazione\\dello switch RF};
\draw[pwr] (dock.east) -| ($(mux.west)+(-6,3)$) -- ($(mux.west)+(0,3)$);
\node[lab,above] at ($(dock.east)+(8,0)$) {5\,V, priorità};
\draw[pwr] (usb.east) -| ($(mux.west)+(-6,-3)$) -- ($(mux.west)+(0,-3)$);
\node[lab,below] at ($(usb.east)+(8,0)$) {5\,V USB};
\draw[pwr] (mux) -- node[lab,above]{5V0} (buck);
\draw[pwr] (buck) -- node[lab,above]{3V3} (load);
\draw[pwr] (69,0) |- node[lab,above,pos=0.75]{5V0} (swv);
\fill[sbred!80] (69,0) circle (0.8);
\draw[ctl] (dock.north) |- (94,21) node[lab,above,pos=0.75,text=sbgreen]{EN\_DOCK: accende il convertitore} -- (buck.north);
\end{tikzpicture}
\caption{Alimentazione. Il dock ha la precedenza sulla USB-C. Il convertitore del 3V3 si accende \textbf{solo} se il dock alza il segnale EN\_DOCK: con la sola USB-C l'MCU resta spenta (capitolo~\ref{ch:power}).}
\label{fig:alim-blocchi}
\end{figure}

\section{Componenti principali}
La colonna \enquote{Rif.} è la sigla del componente sullo schematico e sulla scheda.

\begin{xltabular}{\linewidth}{@{}L{17mm}L{37mm}Y@{}}
\tabhead\hd{Rif.}&\hd{Parte}&\hd{A che cosa serve, per il firmware}\\
\endhead
U301 & STM32L562RET6P & Il microcontrollore (MCU), Cortex-M33, contenitore LQFP64 da 64 pin. La lettera \enquote{P} finale indica la versione predisposta per un alimentatore esterno (SMPS): qui non è usato e i pin VDD12 non sono collegati \fonte{REV-1104 §2.3}.\\
IC204 & Si4464-B1B-FMR & La radio (transceiver sub-GHz): modula, demodula e gestisce i pacchetti. Si comanda via SPI1. Il clock arriva dal TCXO a \SI{32}{\mega\hertz}.\\
Y201 & ECS-TXO-20CSMV4-320-AY & TCXO = oscillatore compensato in temperatura, \SI{32}{\mega\hertz}. Precisione \SI{\pm0.5}{\ppm}, garantito da \SIrange{-30}{85}{\celsius} \fonte{DS-ECS p.\,1}.\\
IC2 & Qorvo QPA9510 & PA = amplificatore di potenza della trasmissione. Si accende quando il transistor Q701 alimenta il suo pin di controllo APC.\\
IC202 & ADI HMC349A & Switch RF: collega l'antenna al ramo TX (porta RF2) o al ramo RX (porta RF1).\\
Q201 & Qorvo TQP3M9036 & LNA = amplificatore a basso rumore della ricezione. Ha un pin per spegnerlo.\\
U701 & 74AUP2G14 & Porta logica invertente: da PC2 ricava il comando opposto per l'LNA (net \net{EN\_RX\_LNA}).\\
IC205 & TI TPS3825-33 & Supervisore: tiene la radio spenta se il 3V3 è troppo basso o se il firmware porta basso PC4.\\
IC4 & TI TPS2121 & Sceglie l'alimentazione a \SI{5}{\volt}: il dock (ingresso IN1) ha la precedenza sulla USB-C (ingresso IN2).\\
IC3 & TI TPSM863257 & Convertitore (buck) da 5V0 a 3V3. Si accende solo con \net{EN\_DOCK} alto.\\
U403 & ST STWD100NYWY3F & Watchdog esterno: resetta l'MCU se il firmware non manda impulsi. Tempo tipico \SI{1.6}{\second}, minimo \SI{1.12}{\second}.\\
U404 & Maxim DS3231M & Orologio (RTC) su I2C2. Mantiene l'ora per qualche ora a scheda spenta grazie a dei supercondensatori.\\
U401 & Infineon FM25V10 & Memoria non volatile FRAM da \SI{1}{Mbit} (\SI{128}{\kilo\byte}) su SPI2. Si scrive senza attese.\\
IC501 & TI TCAN3414 & Transceiver CAN: converte FDCAN1 nei livelli del bus CAN.\\
U402 & WCH CH340E & Convertitore USB$\leftrightarrow$seriale su USART1: console e caricamento del firmware.\\
IC1, IC5 & TI TLV9051 & Amplificatori operazionali delle due sonde di temperatura PT1000 (PA e buck).\\
\bottomrule
\end{xltabular}

\section{Fatti essenziali da ricordare}
\begin{itemize}
\item La radio lavora a \SI{435}{\mega\hertz}, dentro la banda \SIrange{353}{479}{\mega\hertz} del Si4464 \fonte{DS-Si4464 p.\,9}.
\item Un solo pin, \textbf{PC2}, comanda insieme switch RF, PA e LNA \fonte{LOG §41, R-04}. PC2 basso = trasmissione, PC2 alto = ricezione.
\item Il 3V3 esiste solo se il dock alza \net{EN\_DOCK}. Con la sola USB-C l'MCU resta spenta (capitolo~\ref{ch:power}).
\item Il watchdog esterno è \textbf{sempre attivo} in volo: il primo impulso su PC6 va dato entro \SI{1.12}{\second} dall'accensione \fonte{LOG F-12}.
\item Nessun LED è comandato dall'MCU. D203 mostra se la radio è fuori reset, D9 la presenza del 3V3 (solo con il ponticello JP1 chiuso).
\end{itemize}
